
\section{Architecture Comparison and Selection}
The table below gives qualitative engineering hypotheses, \emph{not} results from controlled performance experiments. Actual latency, cost, and accuracy depend on workload, model, external tool latency, and implementation quality.

\begin{table}[htbp]\small\centering
\begin{tabularx}{\linewidth}{@{}p{26mm}p{22mm}p{28mm}Y@{}}\toprule
\textbf{Pattern} & \textbf{Flow control} & \textbf{Latency shape} & \textbf{Main operational difficulty}\\\midrule
Single call & Fixed & One call & Variable response quality\\
Sequential chain & Fixed & Sum of stages & Intermediate error propagation\\
Parallel sectioning & Fixed & Slowest branch + join & Stragglers, reconciliation\\
Parallel voting & Fixed & Slowest vote + judge & Correlated wrong answers\\
Router & Conditional & Router + chosen path & Misclassification, fallback\\
ReAct loop & Model-selected steps & Sum of bounded turns & Unbounded action without harness controls\\
Plan-and-execute & Planned dependency DAG & Plan + critical path & Stale plan and replanning\\
Orchestrator--workers & Dynamic delegation & Supervisor + critical path & Worker scope and synthesis\\
Peer collaboration & Distributed delegation & Negotiation and network dependent & Trust, state and attribution\\
Evaluator--optimizer & Bounded feedback & Sum of revision rounds & Judge failure or non-convergent loop\\\bottomrule
\end{tabularx}
\caption{Workflow, latency shape, and failure mode comparison.}
\label{tab:comparison}
\end{table}

\subsection{Selection algorithm}
\begin{enumerate}
\item Can deterministic code or one schema-constrained model call pass the acceptance test? Prefer it.
\item Are the subtasks predictable and ordered? Use sequential chaining with validation gates.
\item Are subtasks independent and mergeable? Use parallel sectioning. Do not claim voting is independent verification.
\item Is the input type variable but accurately classifiable? Use routing and a typed safe fallback.
\item Must the system iteratively inspect unpredictable external results? Consider a bounded ReAct-style loop.
\item Is a dynamic multi-step dependency plan essential? Consider plan-and-execute.
\item Are the task count and specialization dynamic? Consider orchestrator--workers.
\item Is collaboration needed across independently operated agent services? Consider governed peer agents and A2A.
\item Can quality be objectively measured? Add a bounded evaluator--optimizer loop.
\item Does any step mutate consequential state? Require separately validated authorization and exact human approval, regardless of topology.
\end{enumerate}

\subsection{Measurement model}
For a chain of $n$ stages, a baseline approximation is
\[
 T_{\mathrm{chain}} \approx \sum_{i=1}^{n}(T_i+V_i).
\]
For independent parallel stages,
\[
 T_{\mathrm{parallel}}\approx T_{\mathrm{dispatch}}+\max_i T_i+T_{\mathrm{join}}.
\]
For a loop bounded to $k$ turns,
\[
 T_{\mathrm{agent}}\approx\sum_{j=1}^{k}
 (T_{\mathrm{model},j}+T_{\mathrm{tool},j}+T_{\mathrm{policy},j}).
\]
These formulas omit queueing, retries, rate limits and overlapping operations. Measure p50/p95/p99 latency and \emph{cost per externally verified task}, not just tokens per request.

\section{Applied Reference Architecture: Governed Secure-Code Review}
The example task is to analyze a repository for security risks, correlate SAST, dependency and secret-scanning evidence, and produce an auditable report. It does \emph{not} authorize alert dismissal, default-branch changes, or production deployment. This composition demonstrates routing, read-only parallel collection, a bounded investigation loop, deterministic validation, and optional human approval.

\begin{figure}[htbp]\centering
\resizebox{0.96\linewidth}{!}{%
\begin{tikzpicture}
\node[flowcontrol] (req) at (0,0) {Principal, repo scope, SHA};
\node[flowbox] (route) at (3,0) {Route and plan evidence};
\node[flowbox] (sast) at (6,1.3) {SAST evidence};
\node[flowbox] (deps) at (6,-1.3) {Dependencies and secrets};
\node[flowcontrol] (join) at (9,0) {Deduplicate and correlate};
\node[flowbox] (invest) at (9,-2.1) {Bounded read-only agent};
\node[flowcontrol] (report) at (6,-3.7) {Report verification};
\node[flowreview] (human) at (3,-3.7) {Approve specific write if requested};
\draw[flowline] (req)--(route);
\draw[flowline] (route)--(sast);
\draw[flowline] (route)--(deps);
\draw[flowline] (sast)--(join);
\draw[flowline] (deps)--(join);
\draw[flowline] (join)--(invest);
\draw[flowline] (invest)--(report);
\draw[flowline] (report)--(human);
\end{tikzpicture}}
\caption{Composed secure-code assessment: the default execution path is read-only.}
\label{fig:appsec}
\end{figure}

\subsection{Data and evidence contracts}
\begin{itemize}
\item \textbf{Request:} request ID, repository ID, immutable commit SHA, authenticated principal, classification, allowed tools, deadline, budget and policy version.
\item \textbf{Evidence:} scanner type/version, rule ID, path and range, detection method, retrieval time, immutable reference, severity input and uncertainty.
\item \textbf{Finding:} stable deduplication key, source IDs, affected component, reproduction status, CWE or other taxonomy when supported, remediation proposal and confidence limitations.
\item \textbf{Report:} schema-validated prioritized findings, explicit unknowns, provenance links, verification status, reviewer notes and non-actionable failure records.
\item \textbf{Mutation:} disabled by default. A remediation PR or alert disposition is a new command requiring narrow permissions and separate human approval.
\end{itemize}

\subsection{Pattern alternatives for the same use case}
A small repository with one scanner and deterministic rules may need only a chain. Large repositories with distinct languages can use static fan-out. A heterogeneous workload benefits from routing. Only when follow-up investigation has unpredictable tool paths should an agent loop be introduced. Dynamic delegation is justified by task-specific decomposition; a peer mesh is usually unnecessary unless teams operate independent agent services.

\subsection{Architecture decision record template}
Record context, goals, task distribution, protected assets, stakeholder and regulatory constraints, selected pattern, rejected alternatives, control/permission boundaries, expected latency and cost, failure paths, measurable acceptance tests, evidence from a baseline evaluation, rollout safeguards, and rollback triggers. Document any assumptions that are not yet measured.
